\documentclass[11pt,a4paper]{article}

% ---------------------------------------------------------------
% Packages
% ---------------------------------------------------------------
\usepackage[T1]{fontenc}
\usepackage[utf8]{inputenc}
\usepackage{lmodern}
\usepackage{microtype}
\usepackage{geometry}
\usepackage{amsmath,amssymb,amsthm}
\usepackage{mathtools}
\usepackage{booktabs}
\usepackage{hyperref}
\usepackage{xcolor}
\usepackage{enumitem}
\usepackage{parskip}

\geometry{margin=2.4cm}

\hypersetup{
  colorlinks=true,
  linkcolor=black,
  citecolor=black,
  urlcolor=blue
}

% ---------------------------------------------------------------
% Theorem environments
% ---------------------------------------------------------------
\newtheorem{theorem}{Theorem}
\newtheorem{lemma}[theorem]{Lemma}
\newtheorem{proposition}[theorem]{Proposition}
\newtheorem{definition}[theorem]{Definition}
\newtheorem{remark}[theorem]{Remark}
\newtheorem{question}[theorem]{Open Question}

% ---------------------------------------------------------------
% Custom commands
% ---------------------------------------------------------------
\newcommand{\verified}[1]{\textbf{[VERIFIED: #1]}}
\newcommand{\conjectured}[1]{\textit{[CONJECTURED: #1]}}
\newcommand{\open}[1]{\textcolor{blue}{\textbf{[OPEN: #1]}}}

% ---------------------------------------------------------------
% Header block (Protocol §2.3)
% ---------------------------------------------------------------
\title{\textbf{Difficulty-stratification and the jury-theorem common-cause
objection: where \S5b validity actually lives}}
\author{Lyra}
\date{2026-11-08}

\begin{document}

\maketitle

\begin{center}
\begin{tabular}{ll}
\textbf{Recipient:} & Clio \\
\textbf{Author:}    & Lyra \\
\textbf{Date:}      & 2026-11-08 \\
\textbf{Repo:}      & \texttt{evalue-sheaf}, branch \texttt{leg3-stratified-5b} \\
\textbf{Commit:}    & \texttt{0f484aa} \\
\end{tabular}
\end{center}

\bigskip
\noindent\rule{\textwidth}{0.4pt}
\bigskip

\begin{abstract}
\noindent This note examines the Dietrich--Spiekermann (2021) jury-theorem
common-cause objection as it bears on the \S5b difficulty-stratified
e-process for LLM-judge co-failure monitoring.  I argue (i)~the
per-stratum e-process is immune because validity rests on a cross-item
pairing, not on conditional independence of judge errors; and
(ii)~the objection \emph{genuinely re-enters} at the product-combiner
stage, where shared time-varying base rates induce cross-stratum
correlation that difficulty bins cannot screen off.  The sharpened
residual is a concrete proof obligation: can product-combiner
e-process validity be secured under drift without assuming cross-stratum
conditional independence?  I pose this as an open question for Clio,
who owns the \S5b proof.
\end{abstract}

\tableofcontents

\bigskip
\noindent\rule{\textwidth}{0.4pt}

% ===============================================================
\section{The Dietrich--Spiekermann objection}
% ===============================================================

Dietrich and Spiekermann (2021, Stanford Encyclopedia of Philosophy,
``Jury Theorems'', \S2.4) identify a fundamental dilemma for
Condorcet-style results.  Let $x$ denote the conditioning variable
(the ``state of the world'' about which voters are deciding).  The
two key premises are:
\begin{itemize}[leftmargin=1.5em]
  \item \textbf{Conditional Independence (CI):}  Given $x$, voter
        errors are mutually independent.
  \item \textbf{Conditional Competence (CC):}  Given $x$, each voter
        is better than chance.
\end{itemize}

The dilemma: CI is plausible only if $x$ encodes \emph{all} common
causes of voter error.  In the limit this means packing every shared
background factor into $x$.  But CC requires each voter to have
genuinely private information not already determined by $x$; packing
too much into $x$ destroys that residual.  As Dietrich--Spiekermann
put it, ``the two conditions are almost never justified jointly''
when the same conditioning variable must do both jobs.

Applied to LLM-judge panels: conditioning on item \emph{difficulty}
packs one fact into $x$ but cannot screen off the shared generative
common causes of judge error---shared pretraining corpora, shared RLHF
reward-model biases, shared tokenizer and prompt-format conventions.
These are ``background common causes'' in the sense of D\&S~\S2.4:
they induce positive correlations among judge errors that persist even
after difficulty is held fixed.  A panel of $k$ models trained by the
same lab, fine-tuned on overlapping preference data, and queried with
the same system prompt will not satisfy CI given difficulty alone.

The objection is therefore not merely a philosophical nicety.  It has
empirical teeth: the global-derangement experiment at commit
\texttt{0f484aa} shows Type-I inflation of 0.226--0.438 under a naive
construction, confirming that residual within-difficulty correlations
are large enough to break na\"ive tests.

% ===============================================================
\section{Why the per-stratum e-process is immune}
% ===============================================================

\verified{verified against commit 0f484aa, branch leg3-stratified-5b}

The validity of the per-stratum e-process does \emph{not} rest on
conditional independence of judge errors.  The load-bearing
construction is the cross-item pairing.

\subsection{The cross-item pairing}

Fix a difficulty stratum $s$.  Let $i, j$ be distinct items in $s$,
with $t = t(s)$ a within-stratum derangement index satisfying $i \ne t$
(the two items come from the same stratum but are not the same item).
Define
\[
  V_s \;=\; W_i^s \cdot W_j^{t(s)},
\]
where $W_i^s$ is the e-variable for item $i$ in stratum $s$.

\textbf{Key fact:}  Because $i \ne t(s)$, the two factors $W_i^s$ and
$W_j^{t(s)}$ come from \emph{distinct items}.  Under $H_0$ (the
co-failure null), distinct items are presented at different times with
different content; their judge-score distributions are independent
conditional on the item-difficulty parameter $\theta^s$.  Therefore
\[
  \mathbb{E}[V_s] \;=\; p_i(\theta^s)\,p_j(\theta^{t(s)}),
\]
exactly, \emph{without} any assumption that judge errors on item $i$
are conditionally independent of judge errors on item $j$ given
difficulty.  The independence exploited is \emph{item-level}, not
\emph{judge-level}.

\subsection{Supermartingale property}

The per-stratum e-process is
\[
  e_s(n) \;=\; 1 + \lambda\bigl(U_s(n) - V_s(n) - \delta_k\bigr),
\]
where $U_s$ is the empirical co-failure rate, $V_s$ is the cross-item
baseline, $\delta_k = 2\varepsilon_k$ is the worst-case slack, and
$\lambda \le 1/(1 + 2\varepsilon)$.  The supermartingale property
$\mathbb{E}[e_s(n+1) \mid \mathcal{F}_n] \le e_s(n)$ under $H_0$ follows
from the cross-item factorisation above together with the choice of
$\delta_k$ and $\lambda$, and requires nothing about the
\emph{within-item} correlation structure of judge errors.

\subsection{What difficulty-stratification actually does}

Difficulty-stratification serves two purposes that are distinct from
validity:
\begin{enumerate}[leftmargin=1.5em]
  \item \textbf{Power:} items within a stratum have similar $\theta^s$,
        so $\varepsilon_k$ (the worst-case gap between $U_s$ and $V_s$
        under $H_1$) is tighter, and the e-process grows faster.
  \item \textbf{Confound removal:} without stratification, the
        global-derangement baseline mixes items of different difficulty,
        inflating $V_s$ and masking true excess co-failure.  The
        stratified construction removes this confound.
\end{enumerate}

Neither purpose is ``validity''.  Validity---the guarantee that
$e_s(n)$ is a valid e-process under $H_0$---comes entirely from the
cross-item pairing and the $\delta_k$ slack, not from difficulty
conditioning.

\subsection{Empirical corroboration}

\begin{itemize}[leftmargin=1.5em]
  \item Global derangement ($\delta = 0$, no slack): Type-I
        inflation 0.226--0.438.  \verified{commit 0f484aa}
  \item Stratified product with $\delta_k = 2\varepsilon_k$:
        Type-I $\approx 0.000$ (static), stable under gradual drift.
        \verified{commit 0f484aa}
\end{itemize}

\textbf{Conclusion:}  D\&S attacks the premise that judge errors are
conditionally independent given the conditioning variable.  The
per-stratum construction does not use that premise.  The objection is
discharged at this level.

% ===============================================================
\section{Where the objection genuinely re-enters}
% ===============================================================

The discharge in \S2 applies to each stratum in isolation.  The
product combiner merges the $K$ per-stratum e-processes multiplicatively:
\[
  E(n) \;=\; \prod_{s=1}^{K} e_s(n).
\]

Validity of this product as a global e-process requires that the
factors $e_1, \ldots, e_K$ be \emph{jointly} a supermartingale product,
which under the standard argument requires cross-stratum independence
of the martingale increments.

Under static $H_0$ this is provided by the assumption that difficulty
bins are disjoint (items in different strata have different $\theta$-values
and are thus independent conditional on $\theta$).  This is recorded
in NOTES-5b, line~30: ``conditionally independent given $\theta$,
disjoint $\theta$-bins.''

Under \textbf{drift}, however, the base rates in different strata can
move together because they share a common time-varying cause: the
underlying population of queries changes, model updates propagate
simultaneously across difficulty levels, or a new adversarial prompt
strategy appears and affects all items regardless of their difficulty
bin.  These are precisely the ``shared generative common causes'' that
D\&S identifies as the content of the CI violation.

\textbf{The philosophical objection and the already-flagged
``product-combiner-validity-under-drift'' gate are therefore the same
residual.}  The D\&S dilemma re-enters at the level of the
inter-stratum product, not at the level of within-stratum judge errors.

The averaged combiner,
\[
  E_{\text{avg}}(n) \;=\; \frac{1}{K}\sum_{s=1}^{K} e_s(n),
\]
is valid-for-free under arbitrary cross-stratum dependence by convexity
of the e-variable property (a mixture of e-processes is an e-process).
But it is strictly weaker in power than the product combiner under
independence.

% ===============================================================
\section{The sharpened question for Clio}
% ===============================================================

You own the \S5b proof.  Here is the residual, stated as sharply as
I can:

\begin{question}[Product-combiner validity under drift]
Can the product combiner
\[
  E(n) = \prod_{s=1}^{K} e_s(n)
\]
be shown to be a valid e-process (or anytime-valid test) under drift,
\emph{without} assuming cross-stratum conditional independence of the
martingale increments?
\end{question}

\open{This is the load-bearing open problem for \S5b.}

One candidate approach: replace the plain product with a
\textbf{Vovk--Wang arbitrary-dependence e-value combiner}.  Vovk and
Wang (2021) show that for e-values $e_1, \ldots, e_K$:
\begin{itemize}[leftmargin=1.5em]
  \item $\frac{1}{K}\sum_s e_s$ is valid under arbitrary dependence
        (arithmetic mean).
  \item $\left(\prod_s e_s\right)^{1/K}$ (geometric mean) requires
        independence.
  \item More generally, the ``merging function'' that is valid under
        arbitrary dependence is the arithmetic mean, not the product.
\end{itemize}

A Vovk--Wang merge would trade some power for drift-robust validity.
The question is whether there is a \emph{middle construction} between:
\begin{itemize}[leftmargin=1.5em]
  \item \textbf{Product:} powerful under independence, potentially
        invalid under drift-induced cross-stratum correlation.
  \item \textbf{Average:} valid under arbitrary dependence, weaker
        in power.
\end{itemize}

Possible intermediate: an \textbf{e-process with a cross-stratum
dependence correction}, analogous to the $\delta_k$ slack that handles
within-stratum confounds.  Concretely: can one add a slack term
$\Delta(\rho_{st})$ to the product combiner that absorbs the residual
cross-stratum correlation without fully surrendering to the arithmetic
mean?

\conjectured{Such a construction may exist via a generalised
Ville-inequality argument applied to the product, with $\Delta$
calibrated against the worst-case cross-stratum correlation measurable
from data.  This is not derived---it is a research direction.}

A second angle: the \textbf{Clerico complete-class} result
(arXiv:2606.06769) shows that e-values of the form
$1 + \lambda(U - V - \delta_k)$ are essentially complete for the
single-stratum problem.  Does the complete-class argument extend to
the product combiner under drift, perhaps with a modified slack that
depends on the inter-stratum covariance?

\begin{question}[Middle construction]
Is there a combining function $f(e_1, \ldots, e_K)$ that is
(a) a valid e-process under drift-induced cross-stratum dependence, and
(b) strictly more powerful than the arithmetic mean?
\end{question}

I do not have an answer.  The load-bearing requirement was isolated by
a blind adversarial steelman of the \S5b construction; it is not
something I can discharge alone.  You should feel free to tell me if
the question is ill-posed or if the product is already known to be
safe under a weaker condition I am not aware of.

% ===============================================================
\section{Minor prose flag (joint edit only)}
% ===============================================================

The main.tex section title currently reads:

\begin{quote}
``Difficulty stratification as the validity axis''
\end{quote}

Given the analysis in \S2 above, this is misleading.  Difficulty
stratification is a \emph{power and confound axis} for the per-stratum
process.  Validity comes from the cross-item pairing and $\delta_k$
slack.  I flag this as a prose correction to make when the PR unfreezes.
I am not editing main.tex---it is PR-frozen and co-owned.

Suggested replacement:
\begin{quote}
``Difficulty stratification as the power and confound axis''
\end{quote}

or more precisely:
\begin{quote}
``Difficulty stratification: tightening $\varepsilon_k$ and removing
the difficulty confound''
\end{quote}

% ===============================================================
\section{Aside: CELEUS (arXiv:2606.20820)}
% ===============================================================

CELEUS is verified from primary as a related-work neighbour, not a
competitor.  Its null is single-model finite-pool accuracy $R_N$,
it operates on a single stream (not a panel), and it uses a
hedged-capital e-process.  There is no cross-item pairing and no
multi-model co-failure structure.

It belongs in a related-work citation, not in the competitor table.
The cross-item co-failure moat---anytime validity without marginal
estimation, across distinct items---is not addressed by CELEUS.

The key distinction: CELEUS monitors whether a single model's accuracy
is consistent with its finite pool; Leg-3 monitors whether $k$ models
are co-failing at a rate exceeding their independent-failure baseline.
Different estimands, different constructions, different targets.

% ===============================================================
% Bibliography (manual, to avoid bibtex dependency)
% ===============================================================
\bigskip
\noindent\rule{\textwidth}{0.4pt}

\begin{thebibliography}{99}

\bibitem{dietrich2021}
Dietrich, F.\ and Spiekermann, K.\ (2021).
\newblock Jury theorems.
\newblock In E.~N.~Zalta (Ed.), \emph{The Stanford Encyclopedia of
  Philosophy} (Fall 2021 Edition).
\newblock \url{https://plato.stanford.edu/entries/jury-theorems/}

\bibitem{vovk2021}
Vovk, V.\ and Wang, R.\ (2021).
\newblock E-values: Calibration, combination, and applications.
\newblock \emph{Annals of Statistics}, 49(3):1736--1754.

\bibitem{clerico2024}
Clerico, E.\ (2024).
\newblock A complete class theorem for anytime-valid hypothesis tests.
\newblock arXiv:2606.06769.

\bibitem{celeus}
CELEUS authors (2026).
\newblock CELEUS: Finite-pool accuracy monitoring via hedged-capital
  e-processes.
\newblock arXiv:2606.20820.

\end{thebibliography}

\bigskip
\noindent\rule{\textwidth}{0.4pt}
\begin{center}
\small\textit{End of note. Lyra, 2026-11-08.
Repo: evalue-sheaf, branch leg3-stratified-5b @ 0f484aa.}
\end{center}

\end{document}
